首先感谢王何宇将军三年以来的指导和教学，感谢和我一起选了这门课的同学们，感谢我上过的每一门课的老师和助教，也感谢在工具使用、签到提醒、知识交流和日常共勉中所有帮助过我的人。能落到这个信计2301里读完本科太精彩太牛逼了。

这个人工智能与数学软件怎么说呢？就是我感觉AI这个东西它很牛逼。具体可感地去说它牛逼的话，主要是大语言模型特别牛逼，不过其他的也很牛逼。我记得我第一次看到LLM是在22年年底，也就是口罩刚刚放开的时候。当时在家待了大概一个星期，刷到很多用ChatGPT整活的视频，比如问他一些神人的问题或者引导他出错，或者输出一些戏剧性的东西。当时不太懂这个都是什么东西。然后到上了大学之后，在买手机的时候自己试了一下发现很有意思。如果站在现在去看当年的GPT-3.5，可能觉得它很傻逼，和豆包差不多。但是站在当时来看的话，是很震撼的。一个预训练的东西可以在十秒之内，给一句自然语言的输入，一个不明显输于人类的回答输出。当时觉得非常牛逼，而且泛用性非常非常强。

然后刚上大学的时候吧，我用GPT去写了军训的日记和形策的作业，但是我感觉军训的日记其实写起来，GPT跟百度文库直接搜，或者说直接在范文网站直接搜，也就那么一回事，大差不差。但是如果用来问编程问题就比较牛逼。我记得当时数学分析和高等代数的题目，它基本上能够做到一个六七十的胜率，实际上并没有做得很好，尤其是稍微复杂一点的不定积分，它其实是算不出来的。这里可能暴露出来一个问题，就是GPT-3.5其实并没有真正稳定的推理能力，没训练过的题目能力差一些，本身的推理能力和检索能力有待提升。

比如说我问它梯度下降是什么，它不会说直接把网页上面的东西发给我，可能是把一段教材翻译给我，然后可以随便选出其中的一句话去重新问它。但是也有一些比较特别的点，比如说大一上的时候修了一门课叫药物与健康，是一门通识课，它的题目相当于公开的数据库里面不一定找得到，它做出来正确率就会很差。2023年的GPT还有一个很明显的点是，它的道歉态度非常变态，遇到错就道歉，指出它的问题就马上划归选择另一种对立的答案，也有点类似于现如今的豆包。很抱歉我对你造成了误解，懂的都懂，好吧，懂的都懂，有感觉吧。那时候它最擅长的事情不是推理，是滑跪。你说它错了，它马上说「你是对的，我刚才确实忽略了这一点」，然后给你一个相反的答案。看起来像是在学习，实际上像是在执行某种土下座，然后用户这边如果不仔细看就把这个滑跪合并进去了。

反正当时一开始是觉得很新奇的，也觉得很离谱。也就是这个时候开始，我逐渐开始觉得，自然语言处理，或者说一些比较容易被概括的事情，从翻译、写代码到总结和做报告，这些东西可能还很难说是能力了。我记得当时还有一个很有意思的事情是，信计2023级培养方案里面居然没有四学分的那个大西陈，也就是说，我们这一届是唯一的一个信计里面没有学过大西陈的。我当时觉得这个培养方案改得很诡异，但是大概是在十一月的时候，我感觉GPT-3.5已经很深入地融合到了我的生活之中，然后我开始逐渐领悟到，其实上一个C程的课一点意思都没有。但是这件事情在上了2024的数学软件之后，我才有了更深的感悟。

然后就是2024。当时网站有很多的视频，比如说让GPT去写高考试卷，让它去写考研试卷，或者说让它解决一类的问题。当然我当时并不知道什么是benchmark，什么是baseline。我只知道这些视频还挺有意思的，然后我就逐渐有了兴趣。

24年春夏的时候，也就是24年的上半年，我上一门课叫数据结构基础。这个时候GPT-4已经用得很广泛了，尤其是GPT-4在长文本理解、代码逻辑推理和复杂任务上明显强得多。但是数据结构当时有一个大作业是输入一个数学表达式，含有加减乘除，输出它的求导。这是一个树的遍历的问题，而且这个遍历绝对算不上简单。最后写出来大概是在200到400行的一个代码量是比较正常的。这个代码它其实还是一个单文件编程项目，但是以当时的GPT-4，我记得很清楚，它就算迭代了十几轮也没法给出一个能够跑得通的代码，跑出来会有诸多的问题。当然这个问题比起我们数学人的数据结构，我们的任务好像是输入一个中缀表达式直接把它计算出来，相比于输入一个a乘以a加上a乘以b加上c加上1，输出一个链式法则的求导，需要正确处理链式法则、乘法除法的求导，来说是要简单得多的。

然后呢，GPT-4并不能做得很好。但是从那个时候开始，数学分析和高等代数的题目，在上下文要求不那么长的情况下，已经能够做到几乎全对了。在很多专业和学术的基准上也能达到人类的水平。2023年的时候有一个很热门的话题，是提示词，肯定所有人都有所感受。怎么把自己的背景、要求、格式和评价标准，用一个结构化的、够清楚的方式给大语言模型讲清楚。网上也有很多宣称自己做的提示词惊天动地，或者说别人做的提示词不行，这种风气很厉害，有点类似于现在的人比你的skill不行，我的skill行。

现在回头看，2023年的提示词风潮有点像修仙小说里面的口诀。大家都觉得自己掌握了一段上古咒语，只要把「你是一个资深专家」「请一步一步思考」「输出格式如下」念对了，模型就能开光。后来发现，也不是不行，但是更多时候它只是把香烧得比较整齐。但是武林秘籍这一块还得看我们微分方程数值解。

在我上大学英语三的时候，老师非常明确地跟我们说，我们所有的作文都需要当场完成，因为GPT实在太厉害了。在形势与政策的课上，老师说，我们的报告绝对不允许用AI写。在数据结构的课上，老师说，如果你这个代码是AI写的，你应该把它搞懂。而在接下来2024春夏的数学软件课上，我对很多事情有了新的看法。如果要总结的话，2023年AI的三个主线应该是提示词、作弊鉴定和能力边界的初步探索。

到2024就更加牛逼了。2024年2月的时候，出来了Gemini 1.5 Pro，Gemini一出来就有十多万的token窗口，而且当时已经有一百万token上下文的新闻了，但是客户端没有。如此长的上下文出现之后，AI有能力去处理更大的材料，我们再也不需要在AI输出被截断的时候，输入一个continue或者继续了，它能一口气把很多话讲清楚，也能一口气把很多内容看完。GPT-3.5早期的token上限，如果我没记错的话，应该是比较可怜的，很多时候只有几千到一万多的水平，后面版本才逐渐变长。长上下文从唯物主义上来看，十万的上下文是让AI去阅读文件的必要条件。如果没有十万的上下文，就算支持了文件阅读，AI也无法很好地去阅读这个文件。我们可以让它看一整篇论文之后慢慢讲下来，看完一整组PPT之后决定怎么修改，看完几十页的课程讲义，或者说是很厚的一本书。和原先的输入一大篇文字之后输出一大篇文字，有了很大的变化。

然后是快要期末考的时候，大概是五月到六月份吧，出来了GPT-4o。4o比4多了一个o，这个o应该是omni的意思，也就是全能、多模态。它这个东西能够处理音频、实时音频，还有视觉，也就是图片。这个东西就很牛逼了。比如说小测的时候可以把图片截出来，然后发给它，而不是说先把图片敲成text，把这个分式或者根式表达清楚之后再发给AI让它做题。当然我没干过这种事情。也可以发一个流程图或者表格，也可以发一个电脑的截图，问它按哪个按钮。当然这个东西后面和computer use这种概念有关系，但是我当时不知道。

从这个时候开始，我会有意识地把一些常用的提示词打包起来，放在一个粘贴板应用里面。比如说请把这段话改得去AI化，要求123；请把这段话压缩到原来的一半，要求123；或者请根据这一段教材把XX页到XX页压缩成一份可感可读的期末复习讲义，要求123，等等。我第一次有了管理提示词和存储常用效果的想法。当然后来知道这个东西其实也可以叫Skill。

然后是2024年出现了Sora。Sora不是LLM，但是它能够生成一个非常牛逼的视频或者广告或者短剧。Sora出来的时候反响非常强，因为它生成的视频是整体比较逼真的。不像先前的视频可能还在停留于巨石强森吃石头或者海猫络合物演唱达妹达妹，也就是DeepFake。和先前的那些传统的基于小模型机器学习的尝试完全不一样，它直接从文生视频这个方式入手，有非常强的泛用性，生成的画面也很震撼。

当时是我正式地修读了数学软件这门课程。我们当时的大作业除了几个祖传大作业以外，还有MNIST，相当于是神经网络入门吧，可以这么说。还有一个Mandelbrot集，但是这玩意我觉得确实很数学软件，但是和人工智能没什么关系。还有一些偏向于CPP素养训练和提示词建议的想法，包括老妖对于LLM的使用是一个很宽容的态度，尤其是在数学软件里面基本上是鼓励使用。还有生命游戏，计算π的值，使用线性代数相关的组件完成一个矩阵运算的程序。总而言之，这些都很信计，但是并不是很人工智能，但是比较数学软件。

2024年的5月，我有许多高中同学在刷LeetCode，刷Codeforces。这也是我第一次开始思考这种东西。在三年前，如果我说我高中打过OJ，拿过奖，上大学之后打过ACM，我会编程，我写代码比你快一百倍。确实很牛逼，不得不说就是很牛逼，可能以后也能拿个几百万的工资，但是现在可能并不是这样了。这是我第一次认真地去想，LLM会怎么具体地冲击我们的生活和这个社会。

其实这个时候，我还想到高中时候的一些事情。高中的时候，有些同学在打信息学奥赛，或者自己弄了一个实用小工具，或者搞了一个自己的个人网站。那种东西当时看起来是很酷的。虽然你可能不知道它具体有什么用，但是它本身是有一定技术门槛的。比如说用HTML、CSS、JS搭一个页面，写一点交互，做一个简单的任务。对于零基础的人来说，这肯定不是点一下鼠标就能做出来的，至少要学一会儿，要查点资料，要知道文件怎么组织，浏览器怎么运行，页面怎么调样式，按钮怎么响应。

但是现在有了AI之后，这件事情的感觉变得很不一样。首先，标记语言这种东西感觉能被一脚踢死了。它不是说完全没必要学，而是说这种查表类、积累类、模板类的东西，可以被AI配合得非常彻底。HTML怎么写，CSS怎么居中，JS怎么绑一个事件，Python怎么读一个文件，shell怎么写一个循环，这些东西以前是入门门槛，现在更像是AI随手可以帮你铺好的地板。

更进一步地说，做一些小型项目的门槛大幅度降低了。以前你想搓一个demo，可能要先学框架、学环境、学API、学调试，搞到最后demo还没做出来，人已经烦死了。现在很多时候，只要你有一点程序直觉，知道自己想要什么，有一点设计上的taste，知道这个东西做出来应该长什么样，剩下的就可以让AI以远超先前的效率帮你搭出来。它不一定一次对，也不一定优雅，但是成本和时间已经降到了一个很夸张的程度。

所以我当时第一次很明显地感觉到，编程这个东西的门槛可能真的在变化。以前「我能做出来」本身就很酷，现在「我能做出来」这件事开始变得没那么稀缺了。稀缺的可能变成了：我到底想做什么，我能不能把它做得有用，我有没有taste，我能不能判断这个demo是不是只是一个玩具，还是它真的有一点价值。

而且那个Numerical OJ啊，我用Agent做了一个简单的hack，对不起wcgg啊，对不起。就是那个哈基米哪没路躲，小游戏，我让Agent去看一下它底层是怎么做的，然后直接发了一个最高分的提交上去，然后发现他其实是没有校验和复核的。只需要发一个纯粹的JSON请求包，就可以达到提交成绩的效果。有点类似于我第一次发现学在浙大的quiz，会把答案和题干一起发过来，然后通过脚本可以解析它发的包，进一步可以直接用中间人攻击自动写题目交掉，智慧树也是这样。当然我这个人非常的正直，不会去干这种事情，哈哈哈，就像我不会缺勤，不会翘课，也不会帮人签到，也不会给人抄作业。但是你说放在以前，我这样一个没有CTF经验，也没有web开发经验的人，能做到独立的去调试和解析出一个网站的漏洞，然后再加以利用吗？不管这件事情道德上怎么说，我觉得我技术上都是做不到的。但是Agent的出现大大拉低了这个门槛，甚至纯粹的网页端对话LLM可能也无法帮助我做到这一点。在这之前，我也只是一个会按F12然后装模作样看JSON里面写着什么的普通大学生，当然现在也是。

其实这个时候我仔细想了一个很严重的问题，就是各个群里面都在讨论一个问题：AI会不会取代我们的工作。这个问题本身就很难讨论，因为「AI」是一个非常宽泛的概念，「取代我们的工作」也是一个更加宽泛的概念。这两个宽泛的东西叠在一起，就注定了它会变成一个全民参与、争论不休、谁都能说两句、谁都很难说清楚的问题。这个问题还有点像薛定谔的饭碗。你说它没被取代吧，每天都有人说自己的活已经被AI干了；你说它被取代了吧，大家又还在上班，还在开会，还在填表，还在被ddl追杀。饭碗处在一种既死又没死的叠加态，只有工资到账的时候才会坍缩。

有的人认为，我本科就要去学泛函分析甲，我就要学偏微分方程甲，我要学微分流形、学拓扑，学那些足够抽象、足够高级的课程，让AI无法理解我的课程。有些人认为，既然AI这么屌，迟早要取代，那我不读书了，我去找个岗位当老师吧。还有一些人认为，AI这么屌，那我读什么数学啊，我去做AI吧，感觉还是做AI赚钱啊。还有一些人会问，AI到底可以怎么做数学呢？大家都说AI for Math，AI for Science，那它到底怎么for过去呢？是直接基于大语言模型，还是自己开发一个模型，还是创造更新的理论和管线呢？还有一批人看到以上言论都开始反驳，还有一批人看到什么都不说话。

在这样一个很复杂的话语体系和社区环境里面，我其实也花了很多想法和很多构想去思考自己到底想要做什么样的任务，探索自己的research interest吧。

首先，比起「AI会不会取代我们的工作」这样一个大、宽泛、争论不息的问题，我其实更想讨论的是一个小问题：对于同一件事情，如果我从现在开始学，我能赶得上AI的发展吗？这个问题就比先前那个问题要具体得多，而且答案是比较显然的。那就是不能，至少目前几年看下来不太能。人的学习能力实在太慢了。

虽然你说AI能够写出比如说交换代数的所有习题，它是真的把抽象代数到代数拓扑，到代数几何，再到交换代数所有的课全部都学明白了吗？它是真的领悟了理论的真谛了吗？还是说它只是把交换代数里面的所有习题作为训练素材背下来了呢？不知道。但是我们能看到的是，它确实发展得太快了。

2023年，AI可能还只能写一个一千字以内的单文件程序；2024年，它已经能够在一轮对话里面输出多个文件的完整代码，甚至有机会跑通；到2025年，已经可以用Agent这种形式去完成一个仓库级项目的构建了。又比如2023年，它做数分的题目还会出错，而且甚至还需要我通过清楚的LaTeX把题目讲清楚；2024年，它已经能够处理绝大多数本科问题；到2025年，已经有层出不穷的AI辅助证明、AI辅助数学研究、AI参与重大难题讨论的新闻和爆料了，真假先不说，至少这个方向已经非常热闹了。

不管平时我们在工作的时候怎么抱怨Agent垃圾，怎么夸奖Agent牛逼，在这些情感的、感性的、随口说说的话后面，不得不承认，它的发展实在是太快了。

所以这个时候又回到了一个更核心的问题：我到底有多喜欢数学？我是一个多么纯粹、多么pure的人呢？我觉得我应该不是。所以我可能认命了，也有可能我往后退了，也有可能我跑路了。但是不管用哪种说法，最客观的说法可能是，我对engineering的喜爱超过了对theoretical的喜爱吧。

我不是纯数学人，我可能没有那么pure。我也没有那么那么喜欢去把一本教材或者一篇论文完整啃完。我可能更喜欢engineering。我更喜欢看到自己把东西做出来，喜欢把一个管线跑起来，喜欢看到日志，喜欢看到一个事情按照自己的预期发生，喜欢看到可行的效果，喜欢看到做出来的东西被他人认可，或者切实地加速了什么，简化了我的日常生活，或者优化和改良了工业或者应用上具体存在的问题。

做纯理论当然也很伟大，但是我可能不是那种可以为了一个抽象而精美的理论去纯粹燃烧的人。这个想法其实在我高三填报志愿的时候是比较具体的，但是近来读了两年数学，发现又有点摇摆了。我不止一次地自己想，到底更喜欢做理论还是做工程。但是最后还是觉得更喜欢engineering。

我感觉我也不是突然想明白的，更像是做了几轮人生的梯度下降。高三填志愿的时候给了一个初值，本科读数学读着读着震荡了一会儿，中间还差点掉进一些pure math的局部极小值里，最后好像还是往engineering这个盆地里面滚下去了。至少目前还没发散，局部存在唯一性看起来还行。

后来想想，这里面可能也有一个朴素的no free lunch。AI帮你省掉了写代码的时间，但它会把成本转移到验收、判断、上下文管理和背锅上。它不是把所有困难都消灭了，只是把困难换了一个坐标系。以前难在怎么写，现在难在怎么问、怎么验、怎么知道它有没有在胡说八道。

不过中间会摇摆，说明我们的本科教育质量还不错。哈哈哈哈哈。

然后那个暑假，也就是2024的暑假，我去认真地学习了一些统计机器学习、深度学习，还有一些比较前沿的强化学习和Transformer等等。你说我学了有什么用？它好像确实没有什么用。什么BERT，什么长短期记忆，我操，我现在老早忘了。只知道全部是神经网络里面流动的某种东西而已。通过改神经网络的架构，实现某种功能，然后通过足够的量的训练，就结束了。如果让我现在去默写某个训练流程图，我绝对写不出来。反而是正向传播、反向传播这些比较数学的东西，还记得比较清楚一点。我就记得当时我觉得最速下降法这种东西居然还能带上加速度和速度，感觉创造出来它的人很牛逼。

然后就是到了大二上学期，也就是2024秋冬。2024年9月的时候，GPT出了o1。o1思考时间更长，它的解题能力也更强。这个时候有一个很牛逼的事情发生了，就是对于单文件编程，比如我前面说的那个很复杂的二叉树求导问题，现在o1能够在一个长的时间思考之后，直接给出完整的、已经变过的结果了。远远拉开了与4o的差距。包括数据结构与算法的作业也能够一遍过。至于数学分析和概率论，几乎已经能够达到绝对的正确了，而且在复习的时候能够给我一些直觉上的建议和形象上的说法，而不像是传统的模型那样只能给我一个正确，且也止步于正确的定义或者说教材式的分析，有点像张庆海的教材，但是它却少了一些比较形象直观的东西。但是o1出来之后，这些都变了。它从单纯的能够回答正确的东西，变成了能够持有一个对用户的记忆以及偏好，然后再根据具体的提问来感受你对于这道题的期望是完整、形象，还是冷峻，还是其他的什么样的风格。

我觉得这对大学生活是很重要的。通过这种中间推理过程的学习，让我学习能够强化很多。唯一的不足是ODE，一方面是这门课比较难，另一方面可能计算过程比较复杂，感觉好像AI还是做不太来。甚至当时的AI应该还是存在一个问题，它不会稳定调用一个外部工具，比如说一个Python程序或者一个预设好的tool，它不会去call这个function。这就导致它很难受的一个点是，比如说我让它算一个三乘三矩阵的特征值，然后这个特征值是凑好的三个整数，但是它算不出来，因为中间的多项式和求解，对于自然语言来说，还是太难了。

然后模型的进展也是肉眼可见的很快。o1能够一遍过我们数据结构与算法课程的大作业，也就是那个输入一个中缀表达式输出它的结果，也就是那个关于树和栈的问题。

然后2024年十二月，DeepSeek-V3出来了；到了2025年一月，DeepSeek-R1出来了。DeepSeek牛逼就牛逼，在它首先是国产的，其次它是开源的。在DeepSeek之前应该还有Ollama、通义千问这些开源或者本地运行生态。Kimi和通义千问是什么关系来着？他们是兄弟吗？还是父子？我当时也搞不清楚。哎，反正就那么一回事，但是DeepSeek能力明显强得多，而且在R1出来之后，似乎调用API和修改提示词已经不再是我们能做的全部事情了。我们现在还能够去训练它，去部署它，去微调它，去给它加上一个LoRA的矩阵去改，去修改它的推理流程，或者说把模型直接接入自己的项目系统。

然后当时我在做那啥，数学建模竞赛，虽然那玩意儿后来发现它他妈是个傻逼，但是不管怎么说，大模型能够单独完成几乎整个数学建模任务的时候，人在其中扮演的角色就很微妙了。如果我没有记错的话，几乎就是把一段话发给它，然后再复制回来，再复制过去，再复制回来。相当于我只是一个粘贴板，大多数的思考都是它在做。这可能也是目前大多数学生的现状，可能也是我自己。

我觉得很多人其实都想过，我能不能自己写一个程序，让AI把代码发过来，自动复制它到一个集成编译窗口里面运行它，然后呢，把运行结果丢回给AI，问它合不合理。这个运行结果可能是输出或者报错，然后再把它发过来，然后再跑，再把它发过去，再把它发过来，再跑，直到有一次AI说已经不用修改了，或者说类似的话。其实这个时候已经有harness的雏形了。这个东西实现上应该也是很简单的，毕竟只是在操纵对话框。我当时有这个想法，但是我没把它搞出来，但是后来知道其实是没那么好搞出来的。

然后就是2025年年中，AI编程工具明显升级了，Agent开始变成一个非常非常牛逼的概念。我记得2024的时候，Copilot已经有了，当时的Copilot能干什么来着，反正也是比较傻逼的。它差不多是能够把一行给补得差不多全，然后当时用什么训练的我也不太清楚，它内接了可能是一个数据库吧，反正就不太聪明。然后它可以告诉你一些静态的错误，和一些明显的代码没对上，或者说变量名之前没有提到过之类的，其实还是更靠近一个高级IDE，而不是一个AI。

但是从2025年中，我第一次接触到的是Cursor和GitHub Copilot。当时Copilot我是作为一个插件在VS Code的侧边栏使用的。为什么用它呢？因为它免费，它学生认证免费，而且学生认证的余额很多，等效于一个月卡。更牛逼的是，学生认证可以验证很多个账号，我当时应该拿了四个GitHub Copilot的账号，完全可以满足我的要求。当时也没有五小时限额，只有月限额，一个号用完了就换另一个，用得很开心。可能也是为了推广吧，里面都是最前沿的模型。我也是知道了GitHub Copilot之后才知道了Claude，我以前没听说过Claude。

这里还要感谢Copilot，这是去年我第一次见到现代化的AI Agent，而且还免费了很长一段时间。也感谢whning0315让我有token可以用。不然说实话，很多东西我可能只是听说过，没法真的在自己的作业、代码和项目里面一遍一遍试。

我第一次真正用Agent写作业，大概是2025年的7月左右。那时候有一个数学实践的任务，需要写一个R文件。R文件其实很好写，跟MATLAB差不多一个难度，而且它还是单文件的。但是这件事情给我的震撼不在于R有多难，而在于整个工作流程实在是太奇妙了。

整个流程非常简单，甚至可以把它描述成伪代码。首先，我把题目要求整个截图发给浏览器里面的网页端AI。然后网页端AI给我一段分析。我再和网页端AI说：请把这段分析打包成一个提示词，我会发给Coding Agent。然后我就把这个提示词发给Coding Agent。Coding Agent给我一个输出，我再把这个输出发给网页端AI。网页端AI又给我一个新的提示词，我又发过去；它又给了一个结果，我又发过去；它又给了一个提示词，我又发过去。大概这样持续了五六轮吧，我感觉差不多可以了，就问网页端AI：你觉得可以了吗？它说，我觉得可以了。好，那我就交掉了。

这可能就是所谓的vibe coding吧，只不过我的vibe不是很高级，我的vibe是复制、粘贴、提交、祈祷。更准确一点说，我当时像一个低配版消息队列。网页端AI生产提示词，Coding Agent消费提示词，然后我负责中间的人工Kafka。区别是Kafka不会焦虑，Kafka不会怀疑自己三年本科到底学了什么。

然后我就把它拖到了提交平台上面。这样的日子大概过了一个星期。明眼人这个时候都发现了，我 WHN 在其中扮演的角色难道只是一个 GET/POST 请求吗？我相当于是在用鼠标和键盘操作这台电脑，但是我的所有操作本质上都只是一些网络连接和发包。那我在这个过程中到底学到了什么呢？

虽然数学实验这门课到2026年居然还在教R语言，这件事情本身也有点拉撒，但是我在做这件事情的过程中确实开始意识到：AI Agent不是一个简单的自动补全工具，而是一个可以被我「接线」的工作流节点。网页端AI负责分析题意、组织策略、生成提示词；Coding Agent负责写代码、跑代码、改代码；而我负责在中间传话、判断、确认、最后提交。听起来好像很丢人，像一个人肉网关，但是这就是当时最朴素的Agent工作流。

再比如说，2025年8月的时候，我听说数学分析有一个大项目叫样条，也就是spline。我要去写这个spline，但是我发现我把那个两页的PDF发给网页端AI之后，它也能给出一个长长的提示词。然后提示词发给Agent，Agent自己会给自己运行和验证几轮，然后就做得差不多了。不能说做得特别好，它可能有一些误差分析和结论上的很大问题，但是整体的代码是能写出来的。毕竟样条是一个公开的东西，是很早之前就已经总结得非常标准的技术，并不是什么有探索性或者独创性的工作。这就导致它写起来非常快。

以前是「我本地能跑」，现在是「Agent说它能跑」。这两句话在可信度上其实差不多，都需要验收。有时候它不是没跑通，它是在Agent's imagination里面跑通了。

那我作为这个 GET/POST 机器人，我到底在干嘛？我难道就真的只是守在LLM旁边，然后帮它传递信息吗？那他妈真爆了。

但是后面项目难了之后才会发现，事情并不是这么简单。很多时候需要人去纠偏，需要人去解释具体的结果，需要人判断这个现象到底是因为它本来就应该如此，还是因为代码写错了，还是因为统计口径问错了。这种类似于单元测试，或者说一些直观感受上的事情，我比AI敏锐得多。

比如我可能在三五秒内就能够发现，这个代码应该没错，是它的统计口径错了，导致算出来的结果非常大；也可能发现这个结果不是理论现象，而是代码真的写错了。有些任务看起来只是问一句话，实际上条件数很差。提示词稍微歪一点，结果就爆炸；上下文少一行，方向就偏掉；变量名没读进去，整条管线就开始装死。

以前程序员调试有小黄鸭调试，把代码讲给鸭子听，讲着讲着自己就发现bug了。现在是把报错讲给AI听。区别是小黄鸭不会骗你，AI有时候会一本正经地把你带进沟里。

Agent可能想的是：嗯，这个错误非常明显，它的误差非常大，说明它本来就应该是这样子的，并不是我代码的问题，我应该想办法给出一个结论，用数值分析的语言去解释这个情况为什么会发生。然后它就会开始圆，圆圆圆。但是实际上这个东西本身就是代码出了问题，而不是理论出了问题，那么它圆也是圆不到头的。

所以我后来才意识到，人不是简单地在给Agent传话。人至少要知道什么时候该怀疑结果，什么时候该怀疑代码，什么时候该怀疑统计口径，什么时候该怀疑自己。这个东西听起来很朴素，但是在Agent时代反而变得更重要。

然后，2025年中，我用Agent，我觉得Agent很牛逼。当时我在修读面向对象编程，也就是OOP。OOP区别于前面的编程课或者数据结构课，它的一个明显的区别就是，它是多文件编程。就是OOP可能自己有很多的概念或者理论要遵循，但是最明显的差距就是，它是多文件编程。但是我可以直接把一个文件夹拖进VS Code里，然后跟GitHub Copilot说，你接下来按照这个这么做，把它做好点，不要有AI痕迹，搞快点，等等等等。然后把整个PTA上面的题干发给它，就好了。然后顶多就是我自己再运行跑一下，make一下看看过不过，甚至连GitHub提交，或者说清理硬盘，或者说删缓存，都可以用自然语言去交互了。

24年的时候我修了一门课叫区块链的理论和应用，是陈建海老师的一门通识课。他虽然是通识课，但是比较硬核，然后需要你对Linux有一定的操作水平。我当时不太行，我都是在网页端一边问AI，然后一边在网上搜各种攻略才把它大作业做出来的。如果当时有Agent的话，可能直接一脚踢死。

这个时候我能感觉到，本来我的任务从我能清楚地写出一个代码，我能把代码改成对的，然后自己写一个报告，尽可能快地完成这个任务，这个能力先是变到了我能写出一个完美的提示词，清楚地表达我的意思，然后根据AI生成的结果去有始有终地吸收它输出的结果，然后看懂它给了什么，并且验证结果是否正确。最后到了2025的我能不能清楚地定义任务，拆分模块，让Agent去干这件事情，并且验证结果是否正确，以及怎么把Agent用更有想象力、更有创新力的方法去使用它，等等等等。

然后就是2026，AI Agent很能帮忙整理文件，它能批量处理表格，它能够把自然语言在一个工作区下变成切实可行的任务并且运行它，运行后检查结果，检查是否符合预期，并且继续迭代优化。它能够帮我跑实验，帮我查文件，帮我配环境，Agent能帮我生成代码然后执行，帮我维护项目日志，也可以删除项目日志，可以帮我总结整个课程的资料，然后再把它打包发到GitHub上面，也可以帮我搭建一个自动化的小型的工作流。

它的缺点也逐渐变得明显。Agent会错误操作，Agent会从头造轮子，Agent会搞不清任务的目的，Agent会跑到工作区外面去，Agent会过多地询问安全的行为，Agent会漏过询问危险的行为，Agent会把错误结果包装得像真的。以前大家说LLM是随机鹦鹉，我觉得也没错，只不过这只鹦鹉后来学会了写代码、跑shell、改文件、删数据库，那就有点吓人了。

与其修改代码，它更乐意去自圆其说。Agent会在长的任务里面丢上下文，自己把自己给弄得昏里昏头，不记得自己要干什么，需要很强的权限管理、日志记录和人工验收，不然就会把NumericalOJ的数据库给删掉。所以我觉得在2025年之后，大学生面对AI的一个新能力是，设计一个可控、可验证、可回滚的AI工作流。

有些事故属于概率测度很小，但是损失函数巨大。平时看起来像零测集，真发生一次助教直接进ICU。虽然我还没有遇到过DeepSeek把NumericalOJ的数据库给删掉这种比较恶性的、不可逆的事情，唉，助教太惨了，这种事情太难预防了，但是我感觉进一步收束它的权限，或者说适当地划分一个边界，是非常有必要的事情。

就是说，不管具体在这件事情里面，你做了多少工作，最后负责的时候都要对整个项目有一定的掌控。但是想是这么想的，回头可能又把YOLO开开了。这就是我，太懒了，没办法。我感觉我像个无能的丈夫一样，Agent跟我说搞好了那就是搞好了，它说没搞好那就继续搞，我只负责往里面充钱，然后enter，然后中文汇报，然后继续，然后收到，然后控制c，然后控制v，到点回寝室睡觉。唉，太可怜了，谁来可怜可怜我。

用Agent有时候很像开EVA，同步率高的时候，它一下午给你干完一周的活；同步率低的时候，它在那边说如果你要，我下一步可以顺手再帮你把所有东西都删掉。

而且这个时候，某些research taste或者interest，以及一些宏观上的对产品的感受和目标的规划变得更重要了，具体的落地变成了其中的一部分。虽然Agent还不够很强，但是它已经很强了。

2026年春夏的时候我修读了数值分析。2025一整年我几乎用的都是Gemini，因为Gemini的学生认证不要钱。也是在大概11月12月的时候Gemini 3推出了，我当时觉得Gemini 3非常强。Antigravity也很强，Cursor也很强，但是我还是喜欢在VS Code里面用GitHub Copilot，因为它是免费的，还可以开一堆小号刷余额。

然后就是，学数学到底为我学AI提供了什么？当然我这句话好像假定了学AI是所有人的目标啊，不对啊，不是这样子的。我只是说，我个人学数学为学AI提供了什么。

我第一次从统计学习学到机器学习，再到深度学习和强化学习，以及大语言模型，中间还夹着那些CART、T五、C五点零，还有BERT，还有LSTM。有一些是很具体的公式，有一些是一些流动的东西，我都不知道怎么描述。反正就是软软的东西，不是那么硬硬的东西，SVM或者树肯定是硬的，但是神经网络是一些软的、有弹性的、慢慢在包裹一个解空间的流形。算了，我还是描述不清楚，先不管这个吧，反正就是这个意思。

我第一次自发去学机器学习的时候，其实感觉还是比较吃力的。因为大一刚上完，我只有基本的线性代数基础。虽然线性代数学得不错，微积分也学得不错，数学分析也学得不错，但是当很多误差分析、激活函数的设置，还有一些底层实验都推上来的时候，还是比较吃力的。你会知道它大概在干什么，但是很难真的觉得自己进入了这个语言体系。第一次舍得花钱给一个稳定的梯子，第一次找第三方给GPT充月卡，第一次跑通自己写的ipynb和小实验，第一次ssh到别人的设施上面，第一次accepted，唉，一路走来。但是现在不一样。现在大三了，我学了更多后面的课，学了崔逸凡的那个课，也学了胡贤良的那个课，还有徐翔那个优化，数值代数，数值分析。我感觉我现在再回去看那些训练、loss、过拟合、欠拟合，一直到卷积、旋转、字编码、向量空间映射这些东西，就感觉是一脚踢死，什么路边东西和考驾照科目一一样，哈哈哈哈哈。不是说我全懂了，而是说我好像已经进入了这一套语言体系。对于一类统计和概率的语言，或者说现在2026被夺舍了都说是神经网络的语言，以及这一套体系里面常用的分析方法和直觉，我有了一个更好的把握。

我觉得这些都是回音啊，都是回音。就像是我本科学了数学，然后实习去了量化，科研做的是AI和神经网络，结果最后却莫名其妙地走到了一个AI做量化的地方。感觉是先前的所有铺垫都于此汇聚成了一条浅路啊。「先前的所有铺垫都于此汇聚成了一条浅路。」这句话要加引号，我的过去汇集了，汇集的过去变成了未来，如果我这个人生大概是一个柯西问题的话，那这一步的误差小得挺不错的。

那AI4量化这个方向，数学体现在哪里呢？那当然是链式法则，也就是 partial AI4量化 除以 partial 数学，等于两者分别 partial 数学，哈哈哈，好吧，这个笑话一点都不好笑。但是你懂我意思吧，反正就是这个意思。

我以前学最速下降或者共轭梯度、矩阵分解的时候，可能都觉得这个东西数值代数讲得挺好的，可以用来解线性方程组，对吧？但是后来你会发现，大模型的训练、优化、Adam、SGD、SAG、embedding、LoRA、数值稳定性、病态性，也就是ill-conditioning，全部都是这些东西的影子。其实完全都是homologous的，都是一个妈生出来的。唉，反正就这个意思，你知道吧。

就像是早一点学数值分析的话，可能会觉得数值分析是用具体的代码把数学公式变成脏但是快、而且具体的数值公式，然后再把数值公式变成好用的代码，差不多就是这么一回事。但是在AI时代，在各种工具的帮助下，其实更像是用数学去约束软件，然后再用软件去约束数学。

什么时候该信，什么东西该信，什么时候该检验，什么时候查看误差，误差是多少，大概要花多久，下一步干什么，这些问题都很关键。数学不是只让我会算一个东西，而是让我知道什么东西可能是假的，什么东西应该被验证，什么东西看起来对但其实不可靠。这个东西在AI时代，可能比单纯会手算还重要。

直到2026年的一月份，也就是考完期末考那会儿，我写了一个工作流，全自动地生产竞赛题，生产一种杂交的竞赛题，然后作为数据标注任务发给甲方，赚了几千块钱，然后也真正地感受到了某种可以自动化的，某种有自驱力的和可以轻量批量生产的，轻易批量生产的强大的力量，现在是可以供于每一个人掌握的。

2026春夏学期的时候，我经常在问自己一个问题。包括看前面的内容也可以看到，好像从保存提示词，截图发给AI问它按哪个按钮，到构建一个自主迭代的基于网络聊天的工作流，再到把ls、cd、tree等指令一起作为提示词的一部分发给AI，其实已经摸到了harness、Agent系统提示词、用户提示词、computer use这些内容的概念了。但是我当时为什么没有想出来，构建一个Agent，然后赚大米呢？

我后来想了一下，这个东西其实也是比较历史唯物主义的。Agent出现的一个可能的必要条件是百万上下文，这是我自己的想法。因为二十万上下文的大模型明显不太行，看不下这么多东西。而且对于Agent工作来说，其中是有大量的系统提示词的，是吧。一旦处理不当，短的上下文会导致它记不住任务，从而导致更高的出错率。所以这里可能都是必然，而不是说某个天才突然想出来了Agent。

然后到了2026，那就是人尽皆知到2026了。从25年末至今，Codex还有Claude Code变成了两个最牛逼的Agent。而且是名副其实的最牛逼，一个更偏开放命令行生态，一个更偏闭源商业生态。但是比起这两个Agent，我觉得最牛逼的是王何宇。至少我觉得目前来看啊，会把AI Agent怎么用，以及怎么用来解决问题，怎么正确地去用，放在一个课程里面来讲的人，我好像只见过数学软件了。

另外一方面，Agent变得非常普及。以前我们可能觉得Agent主要可以用来多文件代码，可以让你一句话用一个提示词就能生成一个多文件的代码，一个仓库级的项目，然后自己去检查、运行、检查、迭代、优化，直到它变成一个成品的项目。然后就是，大家都知道了，三月四月出现了OpenClaw，龙虾，这个东西有争议，但是能力边界也非常宽广。再到各种特化的、非特化的、重型的、轻型的，各种Agent的出现，似乎编程和项目构建以及软工开发变成了一个非常非常没有门槛的事情，至少比以前远远要没有门槛。不管是CC98还是B站上面都满地都是各种一眼就能看出来是什么架构的小工具。日常的小工具和日常AI小工具的开发门槛降到了前所未有的低。前无古人地，谁都可以在Agent的辅助下自己搓一个简单的工具出来了。

然后我又很喜欢王何宇说的一句话，如果大家都能把这个东西给做出来的话，要做什么样的东西，做成什么样，问题就会更关键一点。

然后人工智能与数学软件这门课，去年听说要以一个新的课号去开设这门课了，我也是马上拿来了修读。一方面可以刷GPA，一方面也想看看老妖会教什么，as一个已经更新的课程。

2025年的数学软件大作业也挺有意思。一个是手写数字的识别，需要自行训练；另一个是一个偏数值计算的某结构问题，需要自主实现一个数值计算算法，一个比较人工智能，一个比较数学软件。还有一个我给忘了，干什么的来着？反正不是很难，至少2025的数学软件跟2024的数学软件没有什么特别大的差别。都是比较软工和比较OOP的，可以这么说吧，也比较computational，有几个比较trivial有几个比较open，还是挺好玩的，主要是压力小，能学到东西。

我还记得当时准备上数学软件的时候，在CC98搜过经验贴。上面说这门课会教Linux使用和LaTeX的书写，好像还有MATLAB吧，作为短学期课程出现。然后我当时在想，现在有对话式大语言模型了，还需要LaTeX的书写吗？这个这么浅显的背书式标记性语言，标记语言，有什么学习的必要吗？然后果然没讲，说明王何宇应该比我聪明，非常有前瞻性，哈哈哈，这集神了。

然后呢，就是最新的一年，也就是2026的数学软件，改名叫了人工智能与数学软件。我们的作业是写一个Skill，以及祖传的BMP。我不知道这个BMP到底有什么魅力，可能它教会了你怎么处理非文本信息吧。还有一个作业是图片打分的打榜赛。也是通过提供的API去实现一个评分系统。虽然大多数人都做的是纯提示词的，感觉这个东西大有可为啊，它的上限应该非常的高，但是数据集只有一百条的话，很容易过拟合。我之前交作业的时候搞怪交了一个苦力怕上去，居然也给我算到EVAL里面去了。我都想不到那些拿到这一分的人在干什么。那一刻我感觉不是我在测模型，是苦力怕在炸我们。

但是单看身边的人来看，大多数人做图片打榜挑战赛的时候，都是在跟AI Agent说，也就是在一并发送给API的系统提示词里面说，你需要严厉，你需要严厉，你需要严厉，把不对的东西都给我揪出来，就像这样子，许愿式的提示词。有时候AI不好用也不能全怪AI，毕竟garbage in， garbage out。你给它一坨许愿式提示词，它还给你一坨许愿式结果，也算某种守恒。当然这里面也很容易变成AI slop。AI slop最可怕的地方不是它烂，而是它烂得很顺滑，顺滑到你第一眼还以为它写得挺好。看起来很完整，很丝滑，很像那么回事，但仔细一看，里面没有判断，没有结构，没有人的脑子。

那我为什么觉得王何宇这个课非常牛逼呢？我觉得王何宇的课啊，不管是数学软件，还是人工智能与数学软件，还是数值分析，还是数据结构与算法，还是那个计算方法，都很牛逼。我觉得比起一些千百年不换PPT的老师，对着教材讲书上的定理来说，我其实更喜欢这种有前瞻性的、有老师本人思考的、有一些思想上启发的课程。它不只是课上用了最新的工具，而且是把工具应该怎么用变成了课程内容的一部分。三流的老师禁用AI，二流的老师不管AI，一流的老师知道这个东西就在这里，然后基于降低自己花的时间和最大化学生学到的东西去重新设计课。这完全是两种不同的教学态度。当然这句话说得有点狂，但是我确实是这么感觉的。

以前的数学软件教Linux、LaTeX、C++，或者教一些MATLAB，教一些数值计算的库，都很合理。而且放在五年前，那肯定是必然的。这些东西也很实用，对吧？没有人会说会用Linux是不实用的事情。但是现在工具变成了大模型API、多模态、智能体、Claude Code、国产化模型配置。工具变了，数学软件这门课如果完全不变，那才是不合理的。

如果说机械加上智能机械，电气加上智能电气，经管财会土木改成智能班、创新班、卓越班，计算数学改成信息与计算科学很多时候都像是在搞名头的话，我觉得至少数学软件改名叫人工智能与数学软件，是真的在做人工智能，也是真的在对接一些最新的东西，并且把这些东西融入到本科生教育里面。

我对于老妖会在课上面教Claude Code的使用，以及国产化配置，我觉得是非常牛逼而明智的一点。因为今年新出的DeepSeek V4 Preview，DeepSeek V4是一百万Token上下文是吧，而且还是国产的，然后CC也算是半开源的吧，毕竟本身是偏离线的软件，但是Codex不是开源的吗？为什么不用它呢，总之是这个可能啊，可能，我的猜测，刚好绕过了老师课上不好讲外国东西和之前国内的东西性能不够的两个问题，我记得前几年是一直在讲月之暗面的那个Kimi，后来也提到过千问，都是国产的，没毛病啊。当然具体是怎么备课的我也不知道。哎，反正感觉挺巧的。

不过这个课确实有点难以考核。给分给得确实他妈的太好，读计算机的都这么厉害了能不能别来炸鱼了，给计爷跪了；如果考得太难，我们科班同学又搞不出来，那应该怎么考？说不定以后可以搞个时下比较流行的形式，比如说做一个持续两三个月的小项目，或者期末周打一个黑客松。但是这些任务又会因为每个人的模型选择和腰包紧不紧产生区分，对吧？这个东西确实不好搞的，可能进一步的会不会大学这个模式都有点问题？可能阻碍了我们真正高效的学习？上一句话的「我们」有百分之几？前一句话有多大的统计标准误？不知道，可能全给高分但是前几节课把人吓跑确实是个好办法。

我也跟认识的同学聊过。现在这么多人来选数学软件，不是也显得这课设计得好吗？但是为什么没有信计的去选计算方法？后来好像确实没讨论出来。毕竟王何宇自己都没想出来，那我肯定更想不出来。

人工智能到底以后会怎么样呢？现在是什么声音都有。有人说，从金融上看，人工智能如果短期拿不到巨大的突破，资本就会撤资，AI infra会完蛋。有人说，从社会发展规律上看，人工智能取代了太多岗位，同时又没有拿出相应的新岗位，所以它肯定会被禁止。也有人说，从某种臆测的技术上看，Transformer已经走到头了，我们已经开发了Transformer的所有性能，所以以后必须换掉Transformer了。还有人说，Anthropic太傻逼了，迟早把自己玩死，然后GPT一家独大，胜兵必骄，骄兵必败，最后全人类都没有AI可以用了。各种各样的观点非常多啊，我感觉都是傻逼。虽然你让我说它为什么傻逼，我其实一下子说不出来，但是闻着气味就感觉很傻逼。

这可能也是AI做不到的一点吧。我看到傻逼就滑走了，但是AI看到傻逼，可能还要认真分析它到底哪里不对，然后token就用掉了。那么在这一轮对话中，我比AI多省下了四五块钱，我是不是很厉害？对，这么一想我还挺厉害的，确实不赖啊。

遇到大的问题就往上问，遇到小的问题就问网页端AI，遇到清晰的问题就去问Agent，遇到抽象的问题就先不管，然后把token放进老虎机里面，拉摇杆，得到一个似是而非的结果。在这个过程中，好像人的任务从复制到搜索引擎，再把对应的代码复制过来，粘贴到一个地方，开始跑，变成了把Agent的输出结果复制给网页端AI，再把网页端AI的输出结果作为决策发给Agent。确实有点搞，好吧，确实有点搞，但是确实大家都是这么干的，而且这么干效率也确实不低。

当然玩笑归玩笑，我觉得这些讨论真正有意义的地方，不在于预测哪家公司明年死，哪条技术路线一定完蛋，或者哪个岗位三年后全部消失。真正有意义的地方在于，它逼着每个人重新问自己：我到底在学习什么？我到底在生产什么？我到底能不能用这些工具做出以前做不出来的东西？如果所有人都能更快地写代码、更快地做demo、更快地生成报告，那我还有什么东西是值得留下来的？

我觉得还有一个很明显的话题是，在AI变强了之后，工作似乎还没有完全被解决。那么具体是什么样的工作没有被解决呢？抽象地说，就是需要一些宏观看法和统筹能力的问题没有被解决。但是这个说法本身还是太抽象了，并没有讲清楚到底是什么东西。

我之前也和有些人讨论过这个问题，也看别人争论过或者宣扬过各自的观点。站在我目前的工作强度和工作经验上，我的感受是：AI未必能做得像人那么好。这个话听起来很狂，但是我觉得它不是说AI不能写代码、不能跑实验、不能总结报告，而是说对于一个整体的、综合的任务评价来说，有很多东西并不是简单的信息处理。

AI很会给答案，但是它不太会知道这个答案放在真实世界里会不会被人骂。这个能力很微妙，可能叫社会化，也可能叫挨打经验。没有被甲方、老师、助教、同学、评测脚本、服务器权限、磁盘空间和ddl联合殴打过的智能，总感觉还差那么一点人味。

首先是一些经验性质的工程认知。比如说在这台机器上该配什么版本，该设什么项目，该调什么设置；这个模型跑起来大概要几个小时；这个任务如果现在开始跑，晚上能不能出结果；这个bug像不像环境问题；这个结果看起来是模型真的变强了，还是评测哪里漏了；这个项目现在是该继续堆功能，还是该停下来整理数据和日志。这些东西有时候不是几千行上下文能决定的，也不是纯靠信息就能处理的。

其次是taste。也就是我觉得这个任务有没有意义，我觉得这个东西做出来大家会不会用，我觉得这个任务做出来之后能不能发、能不能赚、能不能讲、能不能变成一个更大的东西。AI可以给建议，但是它没有我的处境，也没有我的欲望，更没有我的风险偏好。它不知道我适合做什么样的工作，我喜欢做什么样的工作，我做什么样的工作会有什么样的后果和成果。这些问题AI没法做主，也担不起责任，也成为不了「我」。

有些项目一眼看上去就是「能跑」，但再看一眼就是「只能跑」。这两个字差很多。能跑是程序启动了，只能跑是它除了启动以外好像也没有什么存在的理由。Agent很容易把前者包装成后者，也很容易把后者包装成前者，最后就需要人类来进行开光验收。

第三点其实就是责任。这个问题倒是已经被讨论得比较透彻了。大多数人都认为，需要担责任的工作，比如法务、公务、作为决策手的医生、公安等等，不可能简单地被AI替代。还有一些有仪式性的工作，比如接待、律师、客服等等，也不是把人换成AI就一定更好。还有一些工作，换成AI之后可能更花钱，那就更没什么好说的了。

但是不得不说，在现在的Agent工作环境下，我们坐在电脑前面，把一行指令敲给它，让它去做，然后自己出来在走廊上面走一圈再回去；或者呵斥它哪里做得不对；或者看到它已经变出成果，非常欣喜。有时候我们很清楚我们做了什么，有时候我们只知道Agent说自己做好了。有时候看到限额到了，开始急眼。这些都已经变成每个人的工作日常，也成为了所有牛马的共同话题了。

所以这个时候我到底缺少了什么，或者我应该去争取什么呢？我觉得可能不是去争取「我比AI更会写某一行代码」，也不是去争取「我比AI更会背某个API」，而是去争取几种更上层的东西。

第一是问题定义能力。我能不能把一个模糊的需求，变成一个Agent能做、我能验收、出了问题能回滚的任务。第二是验收能力。我不能只听Agent说它做完了，我要知道它到底改了什么，测试过什么，哪些地方可能是假的，哪些结果只是看起来很像真的。第三是taste。我需要知道什么东西值得做，什么东西不值得做，什么东西是玩具，什么东西有机会变成系统，什么东西做出来是给自己爽一下，什么东西做出来是真的有用。第四是责任感。AI可以帮我跑，帮我写，帮我试，帮我总结，但是最后这个东西能不能交出去，能不能算数，能不能承担后果，还是我自己的事情。

这还远远不是人类补完计划，更多时候是人类补锅计划。AI负责生成，人在后面补测试、补日志、补解释、补责任。也就是说，Agent越强，人越不能只当一个粘贴板。人至少要当产品经理、测试员、架构师、验收人和责任人。

从会写代码，到会验收代码，再到独立完成一个完整的项目，中间到底发生了什么？我觉得这个能力是数学软件这一系列课程一直在教导我们、一直在鼓励我们做的事情。

人写出来的代码看起来对，但是实际上可能有一堆bug。而且如果一个人自己点验自己的代码，可能也没法生成那么多好的bug。到了Agent这个角度，又派生出了一堆新的玩法。比如让Agent自己生成单元测试，然后测自己的代码；让Agent强迫自己说真话；让Agent先自己验收，再把结果交出来；让Agent先设计架构再执行；让它自己去背锅自己的产品；或者让两个Agent相互debate。这些设计都很多，也各有各的美。

但是相信也有很多人被Agent坑过。比如我之前在设计模型的时候，其实是一个比较底层的小问题。当时我在设计一个数据接口，没有安装报错，也没有明显崩掉，但是变量名没有设计正确，设计成了老版本的变量名，导致它始终没有读到这个变量，一直都是空的。进而导致训练出来效果特别差，相比于基线没有提升。那废话啊，他妈的数据变量都没读进去，当然不会有提升。

但是Agent看到这个结果之后，并不会自然地去查里面流淌的东西。它不会主动看管线里面具体流淌了什么东西，它只会看最后流出来的那些东西，看某个统计脚本输出出来的数据是什么样的，然后给这组数据做辩护。它可能会说，这个方法确实无效，我们应该换一个方法。这个时候，如果人输入了一个continue，不管是因为累了，还是因为想先试试吧，还是因为觉得看来这个方法确实不行，在这件事情里面，坐在电脑前面的人就彻底被Agent带偏了。

所以回到这门课，回到人工智能与数学软件，我觉得它最有意思的地方也在这里。它不是简单教你一个软件怎么用，也不是简单教你一个模型怎么调，而是把人、数学、代码、工具、模型、Agent、评价和责任混在一起。它让我意识到，所谓数学软件，可能已经不只是MATLAB、C++、Python、矩阵运算或者BMP文件了。到了现在，它甚至可以是一整套人和AI共同工作的方式：我提出问题，AI拆任务，工具去执行，程序给结果，人来验收，然后再继续迭代。

这个东西很乱，也很危险，也很容易变成老虎机。但是不得不说，它也确实很牛逼。

这个期末大作业我觉得很神圣。之所以神圣，是因为站在一个大三的信计学生的眼里，做完大作业能感觉到自己三年书没一行是白读的，虽然这课不是给大三的学生上的。

这个项目把我前面学过却常常分散在不同课程里的东西，第一次真正压到同一张桌子上：数学分析、线性代数、常微分方程、数值方法、程序设计、算法调试、数据处理、误差意识，最后都不再只是作业题里的局部技巧，而是变成了一个必须彼此配合、彼此负责的完整系统。真正去做月球借力、日心转移、事件搜索、数值积分和结果校验的时候，我很清楚地感受到，所谓「会写代码」和「能把一个复杂问题做成可信结果」之间，其实隔着很长的一段路，而这次大作业让我第一次比较完整地走了这段路。

它也让我意识到，很多看起来抽象、枯燥、甚至当时觉得不知道学来做什么的内容，到了真正要解决一个带有物理约束、数值陷阱和工程细节的问题时，都会一项一项回来，变成你能不能把事情做对的根基。所以我很感谢这次大作业，它不只是让我完成了一份课程任务，更像是让我认真地回头确认了一遍：过去三年的学习，确实已经在我身上留下了能够支撑我解决复杂问题的东西。

一转眼，从大一的时候发现数学软件出分，认识的人互相一问都满绩了，到现在也足足有两年了。我也从海纳苑都没去过的大一新生，变成快要毕业的老狗了。哎，我操，怎么那么快啊？但是本科期间能上这样的本科教育，我感觉也挺满足了。这句话基本不是串的，这篇致谢也是，不过是语音转文字的，这个AI的语音转文字确实有点东西吧，除了人名基本全是对的。

然后呢最后啊，我觉得啊，能够上到这门课是非常有幸的啊，就算他最后给我给分给个3.0，我觉得也是值得的，但是别真的给3.0。务实的课难求啊，真的学到东西啊。我已经稳稳地接住了。
